\chapter{Güven Aralıkları ve Belirsizliğin Yorumlanması}
\index{güven aralığı}
\index{belirsizlik}

\begin{hedefler}
Bu bölümün sonunda okuyucu:
\begin{itemize}
  \item güven aralığının tahmin, standart hata, kritik değer ve hata payı bileşenlerini açıklayabilecek,
  \item güven düzeyi ile gözlenen güven aralığını birbirinden ayırabilecek,
  \item ortalama, oran ve iki grup farkı için uygun güven aralığını kurup yorumlayabilecek,
  \item aralık genişliğini örneklem hacmi, değişkenlik, güven düzeyi ve araştırma tasarımıyla ilişkilendirebilecek,
  \item güven aralığını sıfır değeri ve uygulamada önemli etki sınırlarıyla birlikte değerlendirebilecek,
  \item örneklem hacmini hedeflenen hata payına göre yaklaşık olarak planlayabilecek,
  \item güven aralıklarını Python ve SPSS çıktılarında tanıyıp bilimsel biçimde raporlayabilecektir.
\end{itemize}
\end{hedefler}

\begin{hazirlik}
İki araştırma aynı ortalamayı bulduğu hâlde birinin güven aralığı neden daha
dar olabilir? Örneklem hacmi, değişkenlik ve güven düzeyini düşünerek en az iki
neden yazın.
\end{hazirlik}

\section{PDR vakası: Yalnızlık düzeyi ne kadar hassas tahmin edildi?}

Önceki bölümde birinci sınıf öğrencilerinden olasılıklı yöntemle seçilen 120
kişinin yalnızlık puanı ortalaması 43,6, standart sapması 9,8 bulunmuştu.
Öğrencilerin 36'sı yüksek risk eşiğinin üzerinde puan almış, dolayısıyla risk
oranı yüzde 30 olarak tahmin edilmişti.

Bu sonuçlar iki önemli soruyu açık bırakır:

\begin{enumerate}
  \item Evren ortalaması için 43,6 değerinin çevresinde hangi değerler veriyle makul ölçüde uyumludur?
  \item Evrenin yüksek risk oranı gerçekten yüzde 30'a ne kadar yakındır?
\end{enumerate}

Nokta tahmini tek bir sayı verdiği için bu belirsizliği göstermez. Güven
aralığı, aynı örnekleme sürecinde tahminin ne kadar değişebileceğini dikkate
alarak parametre için makul değerler kümesi oluşturur. Aralık ne kadar darsa
tahmin o kadar hassastır; ancak dar bir aralık, örnekleme veya ölçme sürecindeki
sistematik yanlılıkları ortadan kaldırmaz \cite{casella2002,wasserman2004}.

\begin{researchbox}
Yalnızlık ölçeğinin yüksek risk eşiği klinik açıdan geçerli değilse oran için
çok dar bir güven aralığı elde etmek yanlış sınıflandırmayı düzeltmez. Güven
aralığı yalnız örnekleme belirsizliğini özetler; ölçme aracının geçerliği,
örneklemin temsili ve veri kalitesi ayrıca değerlendirilmelidir.
\end{researchbox}

\section{Nokta tahmininden aralık tahminine}
\index{nokta tahmini}
\index{kritik değer}
\index{hata payı}
\index{standart hata}

Güven aralığının genel yapısı
\[
\text{nokta tahmini}\ \pm\
\text{kritik değer}\times\text{standart hata}
\]
biçimindedir. Kritik değer ile standart hatanın çarpımına \emph{hata payı}
denir. Dolayısıyla
\[
\text{güven aralığı}=\text{tahmin}\pm\text{hata payı}
\]
olarak da yazılabilir.

\begin{table}[htbp]
\centering
\caption{Bir güven aralığının temel bileşenleri.}
\begin{tabular}{@{}p{0.22\textwidth}p{0.30\textwidth}p{0.31\textwidth}@{}}
\toprule
Bileşen & Yanıtladığı soru & Yalnızlık örneği \\
\midrule
Nokta tahmini & Örneklemin en iyi tek tahmini nedir? & $\bar x=43{,}6$ \\
Standart hata & Tahmin örneklemden örnekleme ne kadar değişir? & $s/\sqrt n\approx0{,}89$ \\
Kritik değer & İstenen güven düzeyi kaç standart hata gerektirir? & Yaklaşık $1{,}98$ \\
Hata payı & Tahminin iki yanına ne kadar eklenecek? & Yaklaşık $1{,}77$ puan \\
Alt--üst sınır & Parametre için hangi değerler veriyle uyumludur? & $41{,}83$--$45{,}37$ \\
\bottomrule
\end{tabular}
\end{table}

Standart hata örneklem hacmi arttıkça çoğunlukla küçülür. Kritik değer ise
seçilen güven düzeyine ve kullanılan dağılıma bağlıdır. Güven aralığının
genişliği bu iki unsurun ortak sonucudur.

\section{Ortalama için güven aralığı}
\index{güven aralığı!ortalama}
\index{t dağılımı@$t$ dağılımı}

Evren standart sapması bilinmediğinde ve koşullar uygun olduğunda ortalama
için
\[
\bar x\pm t_{1-\alpha/2,n-1}\frac{s}{\sqrt n}
\]
aralığı kullanılır. Burada $n-1$ serbestlik derecesi, $t$ dağılımından alınan
kritik değeri belirler. Evren standart sapması uygulamada çoğunlukla
bilinmediği için temel araştırmalarda $z$ yerine $t$ aralığı olağan seçimdir.
Örneklem büyüdükçe $t$ dağılımı standart normal dağılıma yaklaşır.

\subsection{Adım adım yalnızlık puanı aralığı}

Yalnızlık verisinde $\bar x=43{,}6$, $s=9{,}8$ ve $n=120$'dir.

\adimbaslik{1. Standart hata}
\[
SE(\bar X)=\frac{9{,}8}{\sqrt{120}}\approx0{,}89.
\]

\adimbaslik{2. Kritik değer}

Yüzde 95 güven düzeyi ve 119 serbestlik derecesi için kritik değer yaklaşık
$t^*=1{,}98$'dir.

\adimbaslik{3. Hata payı}
\[
ME=1{,}98(0{,}89)\approx1{,}77.
\]

\adimbaslik{4. Güven aralığı}
\[
43{,}6\pm1{,}77=[41{,}83;\ 45{,}37].
\]

Veri ve yöntem varsayımları altında evren ortalaması için yaklaşık 41,83 ile
45,37 puan arasındaki değerler yüzde 95 güven aralığıyla uyumludur. Sonuç
``ortalama kesin olarak 43,6'dır'' biçiminde değil, tahmin ve belirsizlik
birlikte verilerek raporlanır.

\section{Güven düzeyinin doğru anlamı}
\index{güven düzeyi}
\index{kapsama oranı}

Yüzde 95 güven, hesaplanan belirli aralığın sabit parametreyi yüzde 95
olasılıkla içerdiği anlamına gelmez. Klasik yaklaşımda parametre sabittir;
örneklem ve ona bağlı aralık değişkendir. Aynı örnekleme ve aralık kurma
işlemi çok kez yinelendiğinde kurulan aralıkların yaklaşık yüzde 95'i gerçek
parametreyi kapsar.

\begin{figure}[htbp]
\centering
\begin{tikzpicture}[alt={Tekrarlı örneklemlerde güven aralıklarının çoğunun gerçek ortalamayı kapsadığını, bir aralığın ise kapsamadığını gösteren çizim},x=0.9cm,y=0.52cm]
  \draw[PrismBlue,very thick] (5,0.25) -- (5,10.55);
  \node[above,font=\scriptsize,text=PrismBlue] at (5,10.55) {gerçek $\mu$};
  \foreach \y/\lo/\hi/\m in {
    1/3.7/5.8/4.75,2/4.1/6.0/5.05,3/3.9/5.5/4.70,4/4.4/6.3/5.35,
    5/3.5/4.8/4.15,6/4.3/5.7/5.00,7/3.8/5.3/4.55,8/4.5/6.1/5.30,
    9/3.6/5.4/4.50,10/5.2/6.7/5.95}
    {\draw[PrismBlue,line width=1.2pt] (\lo,\y)--(\hi,\y); \fill[PrismBlue] (\m,\y) circle (1.7pt);}
  \draw[PrismWarnOrange!90!black,line width=1.6pt] (5.2,10)--(6.7,10);
  \fill[PrismWarnOrange!90!black] (5.95,10) circle (1.9pt);
  \node[right,font=\scriptsize,text=PrismWarnOrange!80!black] at (6.75,10) {kapsamıyor};
\end{tikzpicture}
\caption{Tekrarlı örneklemlerde güven aralıklarının kapsama mantığı.}
\end{figure}

Gözlenen aralık için uygulamada ``evren ortalamasının 41,83 ile 45,37
arasında olduğuna yüzde 95 güveniyoruz'' denebilir. Bu kısa ifade, güvenin
parametreye ilişkin öznel olasılık değil, kullanılan yöntemin uzun dönemli
kapsama başarısına dayandığı bilinerek kullanılmalıdır.

\begin{mistakebox}
Yüzde 95 güven aralığı, bireysel öğrencilerin yüzde 95'inin bu sınırlar içinde
olduğu anlamına gelmez. Aralık evren \emph{ortalamasının} belirsizliğine
ilişkindir. Bireysel puanların yayılımı standart sapma veya yüzdeliklerle
betimlenir ve çoğunlukla çok daha geniştir.
\end{mistakebox}

\section{Güven düzeyi ve aralık genişliği}
\index{güven aralığı!genişlik}

Aynı veri için güven düzeyi yükseltildiğinde daha yüksek kapsama hedeflenir ve
kritik değer büyür. Bu nedenle yüzde 99 güven aralığı yüzde 95 aralığından,
yüzde 95 aralığı da yüzde 90 aralığından geniştir.

\begin{table}[htbp]
\centering
\caption{Büyük örneklemde yaygın güven düzeyleri ve yaklaşık kritik değerler.}
\begin{tabular}{@{}ccc@{}}
\toprule
Güven düzeyi & $\alpha$ & Yaklaşık $z^*$ \\
\midrule
Yüzde 90 & 0,10 & 1,645 \\
Yüzde 95 & 0,05 & 1,960 \\
Yüzde 99 & 0,01 & 2,576 \\
\bottomrule
\end{tabular}
\end{table}

Daha yüksek güven ücretsiz değildir: kapsama olasılığı artarken aralık
genişler ve tahmin daha az keskin görünür. Güven düzeyi sonuçlara bakıldıktan
sonra istenen yorumu elde etmek için değiştirilmemeli; araştırma planında
belirlenmelidir.

\begin{etkinlik}
Bir tahminin standart hatası 2 olsun. Yüzde 90, 95 ve 99 güven düzeylerinde
yaklaşık hata paylarını hesaplayın. Kapsama ile hassasiyet arasındaki
ödünleşimi kendi sözlerinizle açıklayın.
\end{etkinlik}

\section{Ortalama aralığının dayandığı koşullar}
\index{güven aralığı!varsayımlar}
\index{bağımsızlık}

Bir formülü uygulamak, aralığın geçerli olmasını tek başına sağlamaz. Ortalama
için klasik $t$ aralığında şu noktalar değerlendirilmelidir:

\begin{itemize}
  \item Örneklem hedef evrenden uygun bir seçim süreciyle elde edilmiş olmalıdır.
  \item Gözlemler bağımsız olmalı veya sınıf, okul ve tekrarlı ölçüm gibi bağımlılıklar analizde dikkate alınmalıdır.
  \item Küçük örneklemde dağılım güçlü çarpıklık ve ağır aykırı değer göstermemelidir.
  \item Büyük örneklem ortalamanın dağılımını daha kararlı yapar; fakat yanlı örneklemeyi düzeltmez.
  \item Ölçme aracının puanları karşılaştırılabilir ve anlamlı bir nicel ölçek oluşturmalıdır.
\end{itemize}

Örneklem küçük ve dağılım belirgin biçimde çarpıksa dönüştürme, sağlam yöntem
veya yeniden örnekleme seçenekleri düşünülebilir. Karar yalnız bir normallik
testinin $p$-değerine dayandırılmamalıdır.

\begin{assumptionbox}
Önce araştırma tasarımını ve analiz birimini doğrulayın. Aynı danışmandan
hizmet alan öğrenciler ya da aynı öğrencinin tekrarlı ölçümleri bağımsız kabul
edilirse standart hata gereğinden küçük, güven aralığı gereğinden dar olabilir.
\end{assumptionbox}

\section{Bir oran için güven aralığı}
\index{güven aralığı!oran}
\index{Wilson aralığı}
\index{kesin binom aralığı}

Örneklem oranı $\hat p=x/n$ için temel standart hata
\[
SE(\hat p)=\sqrt{\frac{\hat p(1-\hat p)}{n}}
\]
biçimindedir. Büyük örneklemde yaklaşık aralık $\hat p\pm z^*SE(\hat p)$
olarak yazılabilir. Ancak örneklem küçükse veya oran 0 ya da 1'e yakınsa bu
basit aralık kötü kapsama gösterebilir ve $[0,1]$ sınırlarının dışına taşabilir.
Bu durumlarda Wilson aralığı veya kesin binom aralığı tercih edilebilir.

Yalnızlık araştırmasında 120 öğrencinin 36'sı yüksek risk grubundadır:
\[
\hat p=0{,}30.
\]
Wilson yöntemiyle yüzde 95 güven aralığı yaklaşık
\[
[0{,}225;\ 0{,}387]
\]
bulunur. Başka bir deyişle evrendeki yüksek risk oranı için veriyle uyumlu
değerler yaklaşık yüzde 22,5 ile yüzde 38,7 arasındadır. Nokta tahmini yüzde
30 olsa da hizmet planlamasında bu belirsizlik göz önünde bulundurulmalıdır.

\begin{mistakebox}
``Örneklemin yüzde 30'u yüksek risk grubundadır'' gözlenen veriyi betimler.
``Evrende oran tam yüzde 30'dur'' ise belirsizliği yok sayar. Ayrıca güven
aralığı, eşik üzerindeki her öğrencinin gerçekten klinik risk taşıdığını
kanıtlamaz.
\end{mistakebox}

\section{İki grup veya iki zaman arasındaki fark}
\index{güven aralığı!ortalama farkı}
\index{eşleştirilmiş fark}

Karşılaştırma sorusunda ayrı grupların güven aralıklarına bakmak yerine
doğrudan fark için aralık kurulmalıdır. Bağımsız iki ortalamanın farkı genel
olarak
\[
(\bar x_1-\bar x_2)\pm t^*SE(\bar X_1-\bar X_2)
\]
biçimindedir. Eşleştirilmiş ön-test--son-test tasarımında ise her katılımcının
fark puanı $d_i$ hesaplanır ve tek örneklem yaklaşımıyla
\[
\bar d\pm t^*\frac{s_d}{\sqrt n}
\]
aralığı kurulur. Bağımsız ve eşleştirilmiş veriler aynı standart hata
formülünü kullanmaz.

\begin{ornek}[Danışmanlık programı için ortalama fark]
Program grubunun yalnızlık puanı karşılaştırma grubundan ortalama 4,2 puan
düşük olsun. Fark ``program eksi karşılaştırma'' olarak tanımlandığında nokta
tahmini $-4{,}2$, tahmini standart hata 1,7'dir. Büyük örneklem yaklaşımıyla
\[
-4{,}2\pm1{,}96(1{,}7)=[-7{,}53;\ -0{,}87]
\]
elde edilir. Aralığın tamamı sıfırın altındadır; veriler program grubunda daha
düşük ortalamayla uyumludur. Makul farklar yaklaşık 0,9 ile 7,5 puanlık azalma
arasındadır. Fark ters yönde tanımlansaydı işaretler de ters dönerdi.
\end{ornek}

\section{Sıfır değeri, etki büyüklüğü ve uygulamalı önem}
\index{uygulamalı önem}
\index{etki büyüklüğü}

Fark, korelasyon veya regresyon katsayısı için güven aralığı özel bir
referans değerle karşılaştırılır. Ortalama ve oran farkında bu değer 0, oran
biçimindeki ölçülerde çoğunlukla 1'dir. Aralığın sıfırı dışlaması, seçilen
güven düzeyine karşılık gelen iki yönlü testte istatistiksel kanıt bulunduğunu
gösterir.

Ancak araştırma kararı yalnız sıfırı içerip içermemeye indirgenmemelidir.
Uygulamada anlamlı en küçük farkın 3 puan olduğu varsayılsın:

\begin{itemize}
  \item Aralık bütünüyle $-3$'ün altındaysa uygulamada önemli bir azalmayı destekler.
  \item Aralık 0'ı dışlayıp $-3$ ile 0 arasında kalıyorsa fark istatistiksel olarak belirgin, fakat küçük olabilir.
  \item Aralık hem 0'ı hem önemli büyüklükte değerleri içeriyorsa sonuç belirsizdir; ``etki yoktur'' denemez.
  \item Dar bir aralık 0 çevresinde kalıyorsa uygulamada önemli etkinin olmadığı daha güçlü biçimde savunulabilir.
\end{itemize}

\begin{mistakebox}
Güven aralığının 0'ı içermesi grupların eşit olduğunu kanıtlamaz. Veri sıfır
farkıyla uyumludur; fakat geniş bir aralık önemli olumlu ve olumsuz etkileri
de kapsayabilir. Kanıt yokluğu, yokluğun kanıtı değildir.
\end{mistakebox}

\section{Aralık genişliğini belirleyen etkenler}

Yaklaşık toplam genişlik $2\times\text{kritik değer}\times SE$ olduğundan
aralığı şu etkenler belirler:

\begin{itemize}
  \item Güven düzeyi yükseldikçe kritik değer ve aralık genişliği artar.
  \item Örneklem hacmi arttıkça standart hata ve aralık genişliği azalır.
  \item Verideki değişkenlik arttıkça aralık genişler.
  \item Grupların çok dengesiz olması fark tahminini daha az hassas yapabilir.
  \item Kümeleme, ağırlıklandırma ve bağımlı gözlemler standart hata hesabını değiştirir.
  \item Eksik veri nedeniyle etkin örneklem hacminin azalması aralığı genişletebilir.
\end{itemize}

Örneklem hacmi dört katına çıktığında standart hata yaklaşık yarıya düşer.
Fakat aynı kişiden çok sayıda ölçüm almak, aynı sayıda bağımsız kişiden veri
toplamakla eşdeğer değildir. Etkin bilgi miktarı araştırma tasarımına bağlıdır.

\begin{sezgibox}
Güven aralığı bir fotoğrafın netliği gibidir. Daha büyük ve iyi seçilmiş bir
örneklem görüntüyü keskinleştirir; fakat kamera yanlış yöne çevrilmişse yüksek
çözünürlük hedef evren hakkındaki yanlılığı gidermez.
\end{sezgibox}

\section{Hedeflenen hata payına göre örneklem hacmi}
\index{örneklem hacmi!hata payına göre}

Basit rastgele örneklemede, ortalama için beklenen standart sapma $\sigma$ ve
hedef hata payı $E$ biliniyorsa yaklaşık örneklem hacmi
\[
n=\left(\frac{z^*\sigma}{E}\right)^2
\]
ile planlanabilir. Standart sapmanın yaklaşık 10 puan olduğu ve yüzde 95
güvenle hata payının en çok 2 puan olması istendiği durumda
\[
n=\left(\frac{1{,}96(10)}{2}\right)^2=96{,}04
\]
bulunur ve yukarı yuvarlanarak en az 97 tamamlanmış gözlem hedeflenir.

Bir oran için yaklaşık planlama
\[
n=\frac{(z^*)^2p(1-p)}{E^2}
\]
biçimindedir. Ön bilgi yoksa $p=0{,}50$ en koruyucu seçimi verir. Yüzde 95
güven ve en çok 5 yüzde puanlık hata payı için
\[
n=\frac{1{,}96^2(0{,}50)(0{,}50)}{0{,}05^2}=384{,}16
\]
olduğundan en az 385 tamamlanmış yanıt gerekir.

Bu sayılar başlangıçtır. Beklenen yanıtlamama, kümeleme, alt grup analizleri,
sonlu evren düzeltmesi ve bütçe planlamada ayrıca dikkate alınmalıdır. Örneğin
385 tamamlanmış yanıt hedefleniyor ve yüzde 80 yanıt bekleniyorsa yaklaşık
$385/0{,}80=481{,}25$, yani en az 482 kişi davet edilmelidir.

\section{Bootstrap güven aralığı}
\index{bootstrap}
\index{güven aralığı!bootstrap}

Kuramsal standart hata formülü zor olduğunda veya medyan gibi bir istatistiğin
örnekleme dağılımı incelendiğinde bootstrap kullanılabilir. Gözlenen
örneklemden yerine koyarak çok sayıda yeni örneklem seçilir, her birinde hedef
istatistik hesaplanır ve bu değerlerin yüzdelikleri aralık oluşturmak için
kullanılır \cite{efron1993}.

\begin{lstlisting}
import numpy as np

kaygi = np.array([42, 48, 50, 51, 52, 53,
                  54, 55, 57, 60, 62, 85])
rng = np.random.default_rng(209)

bootstrap_medyan = np.array([
    np.median(rng.choice(kaygi, size=len(kaygi), replace=True))
    for _ in range(10000)
])

alt, ust = np.quantile(bootstrap_medyan, [0.025, 0.975])
print("Medyan:", np.median(kaygi))
print("Yuzde 95 bootstrap araligi:", alt, ust)
\end{lstlisting}

Yüzdelik bootstrap aralığı öğretim açısından anlaşılırdır; ancak her durumda
en iyi yöntem değildir. Çok küçük veya temsil gücü zayıf bir örneklemden
yeniden örnekleme yapmak eksik bilgiyi tamamlamaz. Bootstrap da bağımsızlık,
örnekleme tasarımı ve gözlenen örneklemin evreni temsil etmesi gibi koşullara
dayanır.

\section{Kapsama oranını benzetimle görmek}

Aşağıdaki kod, ortalaması 50 ve standart sapması 10 olan bir evrenden 25
kişilik 5.000 örneklem seçer. Her örneklem için yaklaşık yüzde 95 güven aralığı
kurulur ve gerçek ortalama 50'yi kapsayan aralıkların oranı hesaplanır.

\begin{lstlisting}
import numpy as np
from scipy.stats import t

rng = np.random.default_rng(209)
mu = 50
n = 25
tekrar = 5000
kritik = t.ppf(0.975, df=n - 1)
kapsadi = []

for _ in range(tekrar):
    x = rng.normal(loc=mu, scale=10, size=n)
    hata_payi = kritik * x.std(ddof=1) / np.sqrt(n)
    alt, ust = x.mean() - hata_payi, x.mean() + hata_payi
    kapsadi.append(alt <= mu <= ust)

print("Kapsama orani:", np.mean(kapsadi))
\end{lstlisting}

Elde edilen oran her çalıştırmada küçük ölçüde değişmekle birlikte yaklaşık
0,95 olur. Kapsamayan aralıklar hesaplama hatası değildir; yüzde 95 yönteminin
uzun dönemde yaklaşık yüzde 5 aralığı kaçırması beklenir.

\section{SPSS ile güven aralıkları}

\begin{enumerate}
  \item Bir ortalamanın aralığı için \textbf{Analyze $>$ Descriptive Statistics $>$ Explore} yolunu izleyin ve istenen güven düzeyini seçin.
  \item Bir ortalamayı sabit değerle karşılaştırmak için \textbf{Analyze $>$ Compare Means $>$ One-Sample T Test} menüsünü kullanın.
  \item Bağımsız iki grup için \textbf{Independent-Samples T Test}, aynı katılımcıların iki ölçümü için \textbf{Paired-Samples T Test} seçin.
  \item Çıktıda tahminin ve farkın yönünü, güven düzeyini, alt--üst sınırları ve geçerli örneklem sayılarını kontrol edin.
  \item Eksik gözlemler nedeniyle analizlerde farklı katılımcıların kullanılıp kullanılmadığını inceleyin.
\end{enumerate}

SPSS'nin tek örneklem $t$ testi tablosundaki güven aralığı çoğunlukla
ortalamanın değil, örneklem ortalaması ile girilen test değeri arasındaki
farkın aralığıdır. Test değeri 40 ve fark aralığı $[1{,}83;5{,}37]$ ise
ortalama için karşılık gelen sınırlar $[41{,}83;45{,}37]$ olur.

\begin{outputguidebox}
Çıktıyı şu sırayla okuyun: (1) hedef parametre ve farkın yönü, (2) geçerli
$n$, (3) nokta tahmini, (4) standart hata, (5) güven düzeyi ile alt--üst
sınırlar, (6) aralığın sıfır veya uygulamalı önem sınırlarıyla ilişkisi.
Yalnız $p$-değerini okuyup aralığın genişliğini gözden kaçırmayın.
\end{outputguidebox}

\section{Bilimsel raporlama örneği}

``Olasılıklı yöntemle seçilen 120 birinci sınıf öğrencisinin yalnızlık puanı
ortalaması 43,6'dır ($SS=9{,}8$). Ortalama için yüzde 95 güven aralığı
41,83--45,37 puandır. Yüksek risk eşiğinin üzerinde puan alan 36 öğrencinin
örneklem oranı yüzde 30,0; Wilson yüzde 95 güven aralığı yüzde 22,5--38,7'dir.
Aralıklar örnekleme belirsizliğini göstermekte olup örneklemin temsili ve risk
eşiğinin geçerliği ayrıca değerlendirilmiştir.''

Bu rapor nokta tahminlerini, belirsizliği, kullanılan oran aralığı yöntemini
ve yorumun sınırlarını birlikte görünür kılar. Ondalık basamak sayısı ölçme
aracının duyarlılığıyla uyumlu ve metin boyunca tutarlı olmalıdır.

\section{Bir güven aralığını değerlendirirken}

\begin{enumerate}
  \item Aralık hangi parametre için kurulmuştur?
  \item Örnekleme ve araştırma tasarımı kullanılan standart hatayla uyumlu mudur?
  \item Güven düzeyi ve aralık yöntemi belirtilmiş midir?
  \item Aralık yalnız sıfır değeriyle değil, uygulamada önemli değerlerle de karşılaştırılmış mıdır?
  \item Aralığın genişliği güçlü bir karar vermeye yetecek kadar hassas mıdır?
  \item Ölçüm, kapsam, yanıtlamama ve eksik veri yanlılıkları ayrıca tartışılmış mıdır?
\end{enumerate}

\section{Literatür veri setiyle IBM SPSS laboratuvarı}
\index{Student Performance veri seti!güven aralığı}
\index{SPSS!güven aralığı}
\index{güven aralığı!ortalama}
\index{güven aralığı!oran}
\index{Wilson aralığı!SPSS uygulaması}

Bu uygulamada \emph{Student Performance} matematik dosyasından iki güven
aralığı kurulacaktır \cite{cortez2008data,cortezsilva2008}:

\begin{enumerate}
  \item ortalama yıl sonu matematik notu (\texttt{G3}) için yüzde 95 $t$ güven aralığı,
  \item yıl sonu notu en az 10 olan öğrencilerin oranı için yüzde 95 Wilson güven aralığı.
\end{enumerate}

Bir önceki bölümde elde edilen nokta tahminleri böylece örnekleme
belirsizliğiyle birlikte sunulacaktır. Ortalama ve oran farklı parametrelerdir;
bu nedenle aralıklarının hesaplama yöntemleri ve ölçü birimleri de farklıdır.

\begin{researchbox}
Veri setindeki 395 öğrenci iki okuldan gözlenmiştir. Güven aralıkları,
öğrenciler uygun bir örneklem gibi ele alındığında oluşacak örnekleme
değişkenliğini nicelendirir. Okul ve öğrenci seçiminden kaynaklanabilecek
kapsam ya da seçim yanlılığını içermez ve bütün ülkeye otomatik genelleme
sağlamaz.
\end{researchbox}

\subsection{Araştırma görevleri}

\begin{enumerate}
  \item G3 ortalaması için nokta tahminini, standart hatayı ve yüzde 95 $t$ aralığını bulun.
  \item G3 puanı en az 10 olan öğrenciler için 0--1 başarı göstergesi üretin.
  \item Başarı oranı için Wilson yüzde 95 güven aralığını hesaplayın.
  \item İki aralığın genişliğini kendi ölçü birimlerinde yorumlayın.
  \item Güven aralıklarının örnekleme tasarımına ilişkin sınırlılığını açıklayın.
\end{enumerate}

\subsection{SPSS syntax}

\textbf{Explore} yordamı G3 ortalaması için güven aralığını doğrudan verir.
Tek örneklem $t$ testinde test değeri 0 seçildiğinde ``ortalama farkı'' G3
ortalamasına eşit olduğundan, farkın güven aralığı da ortalamanın aralığıdır.

\begin{lstlisting}[language={}]
DATASET ACTIVATE StudentMath.

COMPUTE basarili=$SYSMIS.
IF NOT MISSING(G3) basarili=(G3>=10).
VARIABLE LABELS basarili 'Yıl sonu notu en az 10'.
VALUE LABELS basarili 0 'Başarısız' 1 'Başarılı'.
VARIABLE LEVEL basarili (NOMINAL).
FORMATS basarili (F1.0).
EXECUTE.

EXAMINE VARIABLES=G3
 /PLOT NONE
 /STATISTICS DESCRIPTIVES
 /CINTERVAL 95
 /MISSING LISTWISE.

T-TEST
 /TESTVAL=0
 /MISSING=ANALYSIS
 /VARIABLES=G3
 /CRITERIA=CI(.95).
\end{lstlisting}

Wilson aralığı, örneklem oranı çevresinde simetrik olmak zorunda değildir ve
sınırları 0--1 aralığında tutar. Aşağıdaki syntax geçerli gözlem ve başarılı
öğrenci sayılarını bütün satırlara ekler; ardından Wilson sınırlarını hesaplar.
\texttt{LIST} komutu yalnız ilk satırı gösterir.

\begin{lstlisting}[language={}]
AGGREGATE
 /OUTFILE=* MODE=ADDVARIABLES
 /BREAK=
 /n_gecerli=N(basarili)
 /x_basarili=SUM(basarili).

COMPUTE p_hat=x_basarili/n_gecerli.
COMPUTE z_crit=1.959964.
COMPUTE wilson_merkez=(p_hat+z_crit**2/(2*n_gecerli)) /
 (1+z_crit**2/n_gecerli).
COMPUTE wilson_yari=z_crit*SQRT(
 p_hat*(1-p_hat)/n_gecerli+
 z_crit**2/(4*n_gecerli**2)) /
 (1+z_crit**2/n_gecerli).
COMPUTE wilson_alt=wilson_merkez-wilson_yari.
COMPUTE wilson_ust=wilson_merkez+wilson_yari.
FORMATS p_hat wilson_alt wilson_ust (F6.3).

LIST VARIABLES=n_gecerli x_basarili p_hat
 wilson_alt wilson_ust
 /CASES=FROM 1 TO 1.
\end{lstlisting}

Menüyle ortalama aralığı için \textbf{Analyze $>$ Descriptive Statistics $>$
Explore} yolunda G3, \textbf{Dependent List} alanına taşınır ve
\textbf{Statistics} penceresinde yüzde 95 güven düzeyi korunur. Oran göstergesi
\textbf{Transform $>$ Compute Variable} ile üretilebilir. Wilson hesabı menüde
tek bir standart tablo olarak sunulmadığından syntax ile belgelenmesi işlemi
yeniden üretilebilir kılar.

\subsection{SPSS çıktısının erişilebilir özeti}

\begin{table}[htbp]
\centering
\caption{G3 ortalaması ve başarı oranı için yüzde 95 güven aralıkları.}
\begin{tabular}{@{}p{0.31\textwidth}rrrrr@{}}
\toprule
Hedef parametre & $n$ & Tahmin & SH & Alt sınır & Üst sınır \\
\midrule
Ortalama G3, $t(394)$ & 395 & 10,415 & 0,231 & 9,962 & 10,868 \\
Başarı oranı, Wilson & 395 & 0,671 & 0,024 & 0,623 & 0,715 \\
\bottomrule
\end{tabular}
\end{table}

Başarı göstergesinde 265 öğrenci 1, 130 öğrenci 0 değerini alır. Dolayısıyla
nokta tahmini
\[
\hat p=\frac{265}{395}=0{,}671
\]
ve Wilson yüzde 95 güven aralığı yaklaşık $[0{,}623;0{,}715]$ olur. Tablodaki
0,024 değeri örneklem oranının yaklaşık standart hatasıdır; Wilson sınırları
bu standart hatayı kullanan basit simetrik Wald aralığından farklı bir
düzeltmeyle hesaplanmıştır.

\subsection{Çıktının analizi}

G3 ortalaması 10,415 ve standart hatası 0,231'dir. Yüzde 95 $t$ güven aralığı
$[9{,}962;10{,}868]$ olduğundan yöntem, örnekleme varsayımları altında yaklaşık
9,96 ile 10,87 arasındaki evren ortalamalarıyla uyumludur. Aralığın genişliği
yaklaşık 0,91 puan, hata payı ise 0,45 puandır.

Bu sonuç ``evren ortalamasının bu aralıkta bulunma olasılığı yüzde 95'tir''
biçiminde yorumlanmaz. Sabit olan parametre değil, örneklemden örnekleme
değişen aralıktır. Aynı yöntem çok sayıda benzer örneklemde uygulandığında
oluşturulan aralıkların yaklaşık yüzde 95'inin hedef ortalamayı kapsaması
beklenir.

Başarı oranı yüzde 67,1; Wilson aralığı yüzde 62,3--71,5'tir. Yaklaşık 9,2
yüzde puan genişliğindeki bu aralık, örneklem oranının tek başına vermediği
hassasiyet bilgisini gösterir. Wilson aralığının nokta tahmini çevresinde tam
simetrik olmaması hesaplama hatası değildir.

İki aralığın sayısal genişlikleri doğrudan karşılaştırılmaz: G3 aralığı not
puanı, başarı aralığı oran ya da yüzde puan birimindedir. Her ikisinde de dar
bir aralık yalnız rassal örnekleme belirsizliğinin küçük olduğunu gösterebilir;
ölçüm geçerliğini ve örneklemin hedef evreni temsil ettiğini kanıtlamaz.

\begin{mistakebox}
Tek örneklem $t$ testinde test değeri 0 değilse SPSS'nin verdiği sınırlar
ortalamanın değil, ``ortalama eksi test değeri'' farkının aralığıdır. Ortalama
aralığını bulmak için test değeri alt ve üst sınırlara geri eklenmelidir.
\end{mistakebox}

\begin{outputguidebox}
Çıktıyı şu sırayla okuyun: (1) hedef parametre, (2) geçerli $n$ ve başarı
sayısı, (3) nokta tahmini, (4) kullanılan $t$ veya Wilson yöntemi, (5) güven
düzeyi, (6) alt--üst sınırlar ve aralık genişliği, (7) örnekleme ve ölçme
sınırlılıkları. Güven aralığını olasılık ifadesi veya temsil gücü kanıtı gibi
sunmayın.
\end{outputguidebox}

\subsection{Örnek bulgu paragrafı}

``İki okuldan gözlenen 395 öğrencinin yıl sonu matematik notu ortalama
10,42'dir ($SS=4{,}58$, $SH=0{,}23$, yüzde 95 $t$ GA
$[9{,}96;10{,}87]$). Notu en az 10 olan 265 öğrenci bulunmaktadır; örneklem
başarı oranı yüzde 67,1 ve Wilson yüzde 95 güven aralığı yüzde
62,3--71,5'tir. Aralıklar örnekleme belirsizliğini göstermektedir. Öğrenciler
bütün ülkeyi temsil eden olasılıklı bir tasarımla seçilmediğinden sonuçlar iki
okul bağlamıyla sınırlı yorumlanmıştır.''

\section{Bölüm sonu çözümlü problemler}

\begin{enumerate}
  \item \textbf{Ortalama için güven aralığı.} $n=25$, $\bar x=60$, $s=10$ ve $t^*=2{,}064$ olduğunda yüzde 95 güven aralığını bulun.

  \textbf{Çözüm:} $SE=10/\sqrt{25}=2$ ve hata payı $2{,}064(2)=4{,}128$'dir. Aralık $[55{,}87;64{,}13]$ olur. Bu, bireysel puanların yüzde 95'inin bulunduğu aralık değildir.

  \item \textbf{İki ortalamanın farkı.} Tahmini fark 5, standart hata 2 ve yaklaşık kritik değer 2,01 olsun. Yüzde 95 güven aralığını hesaplayıp yorumlayın.

  \textbf{Çözüm:} $5\pm2{,}01(2)=5\pm4{,}02$; aralık $[0{,}98;9{,}02]$'dir. Veri, birinci grubun evren ortalamasının yaklaşık 1 ile 9 puan daha yüksek olduğu değerlerle uyumludur; sıfır aralıkta değildir.

  \item \textbf{Hassasiyet planlama.} Aynı değişkenlik ve güven düzeyinde hata payını 6'dan 3'e düşürmek için örneklem hacmi yaklaşık kaç katına çıkarılmalıdır?

  \textbf{Çözüm:} Hata payı yaklaşık $1/\sqrt n$ ile değişir. Hata payını yarıya indirmek için $\sqrt n$ iki kat, dolayısıyla $n$ dört kat olmalıdır. Bu ilişki yanlılık ve tasarım sorunlarını çözmez.
\end{enumerate}

\bolumkaynaklari{\cite{casella2002,efron1993,hogg2019,wasserman2004,cortez2008data,cortezsilva2008}}

\section*{Hatırla--Uygula--Yansıt}

\begin{enumerate}
\renewcommand{\labelenumi}{\textbf{\Alph{enumi}.}}
  \item \textbf{Hatırla:} Nokta tahmini, standart hata, kritik değer ve hata payının güven aralığındaki görevlerini açıklayın.
  \item \textbf{Hesapla:} $\bar x=70$, $s=12$, $n=36$ ve $t^*=2{,}03$ için yüzde 95 güven aralığını bulun.
  \item \textbf{Karşılaştır:} Aynı veri için yüzde 90, 95 ve 99 güven aralıklarının genişliklerini sıralayın ve nedenini açıklayın.
  \item \textbf{Yorumla:} Yüzde 95 güven aralığı $[18;26]$ olan bir ortalama için doğru ve yanlış birer yorum yazın.
  \item \textbf{Oran:} 200 öğrencinin 46'sı bir hizmete başvurmuşsa oranı, yaklaşık standart hatayı ve büyük örneklem yüzde 95 aralığını hesaplayın.
  \item \textbf{Fark:} Grup farkı için güven aralığı $[-5{,}4;-0{,}6]$ ise farkın yönünü, istatistiksel kanıtı ve olası büyüklüğünü yorumlayın.
  \item \textbf{Eleştir:} ``Aralık sıfırı içerdiğine göre programın kesinlikle etkisi yoktur'' ifadesini düzeltin.
  \item \textbf{Planla:} $\sigma=15$ ve hedef hata payı 3 puan olduğunda yüzde 95 güven için yaklaşık örneklem hacmini bulun.
  \item \textbf{Uygula:} Python benzetimindeki tekrar sayısını ve örneklem hacmini değiştirerek kapsama oranı ile ortalama aralık genişliğini karşılaştırın.
  \item \textbf{Yansıt:} Dar bir güven aralığının neden geçersiz bir ölçme aracını veya yanlı bir örneklemi telafi edemeyeceğini PDR bağlamında tartışın.
  \item \textbf{Raporla:} Bir ortalama ve bir oran için tahmin, yüzde 95 güven aralığı, örneklem hacmi ve yorum sınırlarını içeren kısa bir bulgu paragrafı yazın.
\end{enumerate}